\section{The five numerical clauses for the canonical family table}

% BEGIN N78T1
\paragraph{N78T1. Reading the entries on one hub.}
Fix the original size datum $L$ and canonical hub $X$ in a family
setup. Let $\ell=\ell_i$ be the prime at index $i$, $n_\lambda$
the local degree of $E$ at $\lambda$, and $m,r,t$ the ranks in
$D.R$. Put
\[
 U=x^\epsilon,\quad b=\exp U,\quad S_c=L.\mathrm{Sc},\quad
 \eta=L.\mathrm{etaSz},\quad \eta_3=\eta+t+2,\quad h=b+S_c+\eta_3 .
\]
Here $x>1$, $\epsilon>0$, $S_c,\eta>0$, so $U,b,h\ge1$.
Write $w''=b'/(xn)$ for the given lower weight,
where $b'=xb/\exp(x^{\epsilon/3})$, and put
\[
 T=(t+1)/2,\quad
 w_B=\max\{w''+\eta,\ \sup_{v\in\mathrm{TT}}b_v+b'/x+\eta\},
 \qquad w_P=w_B+T,\qquad G_P=w_P+S_c+\eta_3 .
\]
In particular $w'',w_P,G_P$ are nonnegative. In the canonical
table $\mathrm{tabX}$ the band and exceptional set are $X.\mathrm{band}$
and $X.\mathrm{exc}$. Its length $d_{i\lambda}$ is the maximum of
$X.\mathrm{len0}(i,\pi,\lambda)$ over the finite set $Z_{\rm good}(i)$,
with value zero when that set is empty. In each row below the
canonical data of DFD8 assign the same entry to every $\pi$. Thus
the maximum is that entry on a nonempty set. A nonnegative upper
bound for it also covers the empty set.

Use the original row constants $R$ and a common real constant $C$
with $R.\mathrm{Le}(C)$. In particular $C\ge1$, $C_A,C_B\ge0$,
$C_2\ge1$, and
\[
 \begin{gathered}
 3C_2+1,\quad 2(3C_A+1),\quad 6C_B+1,\quad
 K_1,\quad K_2+1,\quad K_3\ \le C,\\
 K_1=C_B+48,\qquad K_2=6C_B+12+6C_C,\qquad
 K_3=3C_AC_B+12C_A+9 .
 \end{gathered}
\]
For the actual family $C_B=320$, $C_C=25$, $C_A=23$.
Then $(K_1,K_2+1,K_3)=(368,2083,22365)$, $2(3C_A+1)=140$,
and $6C_B+1=1921$. Thus any $C\ge\max(3C_2+1,22365)$ satisfies
the entire row-constant condition. The source's $C_s$ has these
two lower bounds before $x,i,\lambda$ or the exceptional set is
chosen. This observation will allow all five clauses to use $C_s$.
% END N78T1

% BEGIN N78T2
\paragraph{N78T2. Row A, including the doubled split length.}
For a split prime $\ell\ge5$ put $a=\lfloor h\rfloor$ and
\[
 Y=1+h/\ell+\log_\ell h,\quad e_P=CY,\quad
 e_A=\big\lceil C_A(2+2a/\ell+\log_\ell(2a+1))\big\rceil .
\]
Since $0\le a\le h$ and $h\ge1$, we have $2a+1\le \ell h$.
Consequently $\log_\ell(2a+1)\le1+\log_\ell h$. The expression
in parentheses in $e_A$ is between zero and
$3+2h/\ell+\log_\ell h\le3Y$. The rounding error is at most one,
and $Y\ge1$, so $e_A\le(3C_A+1)Y$. N78T1 gives $2e_A\le e_P$.

For every $\lambda$ the split declaration makes the common good-point
entry $(2n_\lambda m+1)e_A$ on its column set and zero elsewhere.
Thus $2d_{i\lambda}\le(2n_\lambda m+1)e_P$ in either case, also
when the good-point set is empty. Since $\ell$ is nonzero the set
of primes of $E$ over $\ell$ is finite. Summing proves
\[
       2\sum_{\lambda\mid\ell}d_{i\lambda}
          \le \sum_{\lambda\mid\ell}(2n_\lambda m+1)e_P .
\]
This gives row A at every original twisted index with
$\ell\le\ell_*$ or $\ell\in X.\mathrm{exc}$ and with $\ell$ outside
the two distinguished primes.
% END N78T2

% BEGIN N78T3
\paragraph{N78T3. Row B at small inert primes of large order.}
In the table an inert prime may have either the signal or the ordinary
inert class. Denote by $o_\ell$ the least positive exponent giving
$\ell^{o_\ell}\equiv\pm1\bmod 2^{e+2}$ when such an exponent exists,
with infimum zero for the empty subset of $\mathbb N$.
For row B the guards are $\ell<\ell_A$ and $4L_w<o_\ell$, at a
twisted index outside the distinguished primes. Put
\[
 j=\max(1,\lceil\log_{\ell-1}\ell_A\rceil_{\mathbb N}),\qquad
 B_\lambda=(n_\lambda+1)r
 \left(\frac{h}{(\ell-1)^{j+1}}+(G_P+1)(j+1)+U\right).
\]
These logarithms use the source convention
$\log_b y=\log y/\log b$ with division by zero returning zero
if the base is one. The canonical choice $\operatorname{ratIn}(u,v)$
is a rational strictly between $u$ and $v$ when $u<v$, and zero
otherwise. Here $\ell\ge2$, $r>0$ by the rank condition $m\ge2$,
$U\ge1$ and every other summand in parentheses is nonnegative.
Thus $B_\lambda>0$, so the interval
\[
      \left(3C_2 B_\lambda,\ (3C_2+1)B_\lambda\right)
\]
is nonempty. The two strict row guards select
$\mathrm{lenRowB}$, whose value is $\operatorname{ratIn}$ of those
endpoints. The inert declaration sets every good-point entry to this
value. Its upper endpoint is at most $CB_\lambda$ by N78T1.
The empty maximum is zero and satisfies the same inequality.
This proves $d_{i\lambda}\le CB_\lambda$ for every $\lambda\mid\ell$.
% END N78T3

% BEGIN N78T4
\paragraph{N78T4. The medium inert row and its quadratic coefficient.}
Fix an inert $\ell$ outside row B. Write $I_\lambda=1$ if
$X.\mathrm{band}(i,\lambda)$ and $I_\lambda=0$ otherwise. Define
\[
 \begin{aligned}
 f_\lambda&=(w_P+\eta_3)/\ell+S_c+U+\eta_3 I_\lambda,
 &Y&=1+h/\ell+\log_\ell h,\\
 H_\lambda&=(n_\lambda+1)r\,h/\ell^2,
 &F_\lambda&=(n_\lambda+1)r f_\lambda,\qquad
 M_\lambda=n_\lambda mY .
 \end{aligned}
\]
All three quantities $H_\lambda,F_\lambda,M_\lambda$ are
nonnegative, using $\ell\ge2$, $h\ge1$ and N78T1. The inert
complementary row assigns each good point the rational selected
between $K_1H_\lambda+K_2F_\lambda+K_3M_\lambda$ and
$K_1H_\lambda+(K_2+1)F_\lambda+K_3M_\lambda$, with the same
zero convention for an empty interval as in N78T3. Since
$K_1,K_2+1,K_3\le C$ and $C\ge1$, the upper endpoint is at most
\[
          C H_\lambda+C F_\lambda+C^2 M_\lambda .                 \tag{*}
\]
Both the default zero and the empty good-point maximum satisfy this
nonnegative bound. A medium inert prime satisfies
$\ell_A\le\ell$, contradicting row B's first strict inequality.
Set $e_P=CY$. Combining the first two summands in (*) gives
\[
 d_{i\lambda}\le C(n_\lambda+1)r(h/\ell^2+f_\lambda)
                    +C n_\lambda m e_P .
\]
This is the medium inert bound at the prescribed twisted indices
with $\ell\le\ell_*$ or $\ell\in X.\mathrm{exc}$ outside the
distinguished primes. The second factor $C$ in the last term is
already inside $e_P$.
% END N78T4

% BEGIN N78T5
\paragraph{N78T5. The small inert row at the order boundary.}
Suppose $\ell<\ell_A$ and $o_\ell\le4L_w$ at a twisted inert index
outside the distinguished primes. The second inequality excludes the
second strict guard of row B, including the equality case. Hence the
same canonical complementary rational entry occurs as in N78T4, with
the band indicator for this very $(i,\lambda)$. Keep the first two
terms of (*) separate to obtain
\[
 d_{i\lambda}\le C(n_\lambda+1)r h/\ell^2+
                  C(n_\lambda+1)r f_\lambda+C n_\lambda m e_P .
\]
This is the small-order row for each $\lambda\mid\ell$. Neither the
medium-prime cutoff nor membership in the exceptional set is needed
in this clause.
% END N78T5

% BEGIN N78T6
\paragraph{N78T6. From a non-thick letter to the band estimate.}
Assume the setup's parameter inequality $S_c+\eta_3\le b$.
For a twisted index and a marked pair $(i,\lambda)$, the canonical
definition of the band gives $\ell_i\ge11$, $\lambda\mid\ell_i$,
a good point, and word data witnessing the mark. Any rational prime
$\ell$ below $\lambda$ equals $\ell_i$: two distinct rational primes
are coprime, and their integral B\'ezout identity would put $1$ in
$\lambda$. Use this same $\ell$ in what follows.

Write $b_0,w_0,S_0,\eta_0$ for the word parameters. A letter is
thick when its size lies in $[b_0,b_0+\eta_0]$. The band witness
selects a non-thick letter of size $z$ adjacent to a thick one and gives
\[
 |\ell z-b_0|
    \le C_B(\ell\eta_0+\ell S_0+b_0/\ell+z+\ell g),\qquad
 g=\log_\ell(2h_0).
\]
The word height $h_0$ is $b_0+\eta_0+S_0$ if some letter is thick,
and $w_0+S_0$ otherwise. Since $20w_0\le b_0$, in either branch
$h_0\le b_0+\eta_0+S_0\le h$. Also $h=b+S_c+\eta_3\le 2b$.
Therefore $2h_0\le4e^U$, and $g\le U$: for $h_0>0$ this follows
from $\log4+U\le U(\log4+1)\le U\log\ell$, since $U\ge1$ and
$4e<11\le\ell$. If $h_0=0$, the total real logarithm has value zero.

The word box gives $s$ with $w''\le s\le w_B$ and $|z-s|\le T$,
and gives $b-T\le b_0\le b-T+1$. In particular $|b-b_0|\le T$
since $T\ge1/2$. The base bounds are $\eta_0\le\eta_3$ and
$S_0\le S_c$. Set
\[
      E=\ell\eta_3+\ell S_c+b/\ell+s+\ell U .
\]
Then $E\ge\ell\eta_3\ge0$, since $s\ge w''\ge0$. The expression
multiplied by $C_B$ in the word estimate is at most
$E+T+1/\ell\le E+\ell\eta_3\le2E$;
here $2T+1\le\eta_3$ and $\ell\ge11$. The difference between
$\ell s-b$ and $\ell z-b_0$ is at most $\ell T+T\le\ell\eta_3$.
The triangle inequality and $C\ge6C_B+1\ge2C_B+1$ now give
\[
  |\ell s-b|\le(2C_B+1)E\le CE .
\]
This proves the band clause with both endpoints on $s$ and every
prime allowed in its quantifier.

Together N78T2--N78T6 give $\mathrm{RowNum}(L,X,C_s)$ on the
original table $\mathrm{tabX}(L,X)$. This is the $\mathrm{num}$
field of the family conclusion. The good-point maximum, the band
witnesses and the exceptional set all belong to the hub chosen
for that index.
% END N78T6
